\subsection{Comparison with \texorpdfstring{$\Lambda$}{Λ}CDM Cosmology}
  \label{subsec:comparison-with-lambdacdm-cosmology}

  $\Lambda$CDM introduces cold dark matter, dark energy, and inflation as effective
  postulates~\cite{peebles1993principles,planck2020results}.
  In Cosmochrony, cosmic expansion follows from the monotone growth of projected capacity,
  with $H(t) = \dot{\chi}/\chi$ and $H_0 \sim c/\chi(t_0)$.
  Dark energy is reinterpreted as the large-scale growth of projected capacity;
  cosmic acceleration reflects the cumulative advance of projective ordering over
  cosmological timescales.
  The coincidence problem and the Hubble tension may admit interpretations
  in terms of epoch-dependent projective ordering structure rather than
  new fundamental constituents.
  Low-$\ell$ CMB suppression is explored as a structural consequence of
  projective admissibility constraints rather than a statistical fluctuation.

  Table~\ref{tab:comparison_cosmochrony} summarizes the conceptual
  differences between Cosmochrony, $\Lambda$CDM, and Loop Quantum Gravity.
